% RESEARCH_DRAFT_SHA256: 82ed26378b8743cd3cad4b101af828f0bfd6c663cc5a34e3fe20d700aa30caf1
% Layout and source-reference wrappers expanded for the web renderer.
\documentclass[11pt]{article}
\usepackage[margin=25mm]{geometry}
\usepackage{fontspec,unicode-math}
\setmainfont{Latin Modern Roman}
\setmathfont{Latin Modern Math}
\usepackage{amsmath,amssymb,mathtools,xurl,hyperref}
\hypersetup{colorlinks=true,urlcolor=blue,linkcolor=blue}
\setcounter{secnumdepth}{0}
\setlength{\parindent}{0pt}
\setlength{\parskip}{5pt}
\setlength{\emergencystretch}{3em}
\newcommand{\sref}[2]{\hyperlink{source:#1:#2}{[S#2]}}
\newcommand{\cataloguescope}[1]{\paragraph{Source record.}#1}
\newcommand{\floor}[1]{\left\lfloor#1\right\rfloor}
\begin{document}
\section*{Research-attempt provenance}
This standalone research attempt retains the complete original definition body below,
including its historical source and preparation notes. Those notes describe the
earlier catalogue preparation; the new attempt and its verification are identified
separately after that body. No original ID, catalogue status, or source statement is changed.

% UnsolvedMath ID 20002175; source catalogue code AIM-INFRASTRUCTURE-0075
\section{A safe publication-first launch contract}\label{20002175}
{\small UnsolvedMath ID 20002175 · Computer Science}
\subsection{Definitions and mathematical statement}

\paragraph{Shared notation.}Standard number systems are $\mathbb N,\mathbb Z,\mathbb Q,\mathbb R,\mathbb C$. Specialized notation and index ranges follow the cited source and the definitions below.


An XML publication element is a structural node in publication configuration. Content source and publication configuration serve distinct purposes: one describes what the document contains, while the other can specify how it is produced. The exact placement of an element is governed by the system's schema.
\paragraph{Record scope and qualifications.}
The catalogue title is retained as metadata. Named certificates, contracts, bounds or frameworks in that title are not asserted as results of the original source extract.
The isolated element name is preserved literally. It does not specify a complete XML document or establish that such a root is accepted by current software.

Catalogue screening limitation: PreTeXt software design, documentation, or publication feature request; no mathematical problem is specified by the supplied passage.

The source snapshot labels this record partially solved; this is catalogue metadata rather than independent verification of a mathematical result or the current completion of a project.
The source text below is rendered literally where necessary; damaged formulas and composite extracts are not silently repaired into a new mathematical claim.
\paragraph{Statement recorded in the supplied source.}
Start a document with <publication> . \href{https://github.com/PreTeXtBook/community/wiki/Moderated-Problem-Session}{[S2]}
\paragraph{Source status.}
The supplied catalogue labels this source record partially solved; that metadata was checked against the supplied snapshot on 2026-10-11. The label is preserved for provenance and is not an independent verification of present open status or mathematical correctness.
\subsection{Short English statement}
Consider starting a document workflow with publication configuration.
\subsection{Sources}
\begin{itemize}
\item[S1] Ulam AI and the UnsolvedMath Contributors, UnsolvedMath: A Curated Collection of Open Mathematics Problems, record 20002175 (AIM-INFRASTRUCTURE-0075). Catalogue curation and metadata: CC BY 4.0. The source extract is retained; definitions, English summaries and qualifications are editorial additions.
\url{https://huggingface.co/datasets/ulamai/UnsolvedMath}.
\item[S2] AIM Workshop Problem Lists, source linked by catalogue record AIM-INFRASTRUCTURE-0075.
\url{https://github.com/PreTeXtBook/community/wiki/Moderated-Problem-Session}.
\end{itemize}
\paragraph{Preparation and review.}
The original source and UnsolvedMath contributors retain attribution for the source extract. Codex (GPT-6; exact serving revision not exposed) assisted with definitions, English summaries, scope notes and formatting. No new human mathematical authorship or independently verified solution is asserted. Permanent source ID 20002175 and code AIM-INFRASTRUCTURE-0075 checked against the complete supplied registry and current repository. Distinct upstream records, including repeated questions and broad or damaged extracts, are retained for complete coverage of the source CS category. A repeated record is not a new independent conjecture. Complete mathematical-equivalence and current-literature review has not been performed. Preparation and final review follow contribution prompts C and D, with source gaps disclosed. Structural and rendering checks do not establish a new theorem or open status.
\label{end:20002175}

\section{Literature and context check}
The primary community-session source was consulted. The record is a broad software-design request, and no mathematical theorem follows from its editorial catalogue title. The finite model below is explicitly auxiliary; it does not certify an actual implementation or current feature status.
\section{Research attempt: A publication-first wrapper can create an import cycle}
Represent a publication entry $p$ and a content document $b$ as nodes of a mandatory import graph. If $p$ imports $b$ and $b$ imports $p$ to recover its configuration, a naive recursive loader alternates $p,b,p,b,\ldots$ and never finishes. Simply choosing publication as the first node does not remove this cycle.

For a finite graph under one-time import semantics, a loader can distinguish unvisited, active, and completed nodes. On visiting an unvisited node, mark it active and visit its dependencies; then mark it completed. Encountering an active node exhibits a directed cycle along the current recursion stack. Encountering a completed node reuses its result. Each node and edge is processed at most once apart from constant-time state checks, so the procedure terminates in $O(|V|+|E|)$ time and either reports a cycle or yields a dependency order. In the latter case, each node's dependencies are completed before the node, so deterministic evaluation of nodes using only those results is well-defined.

This proof treats cycles as invalid. A language could instead assign fixed-point semantics to some cycles, but that is an additional specification and requires an existence and uniqueness analysis; the graph algorithm does not provide it. For instance the mutually dependent numeric equations $p=b$ and $b=p$ have many solutions, while $p=b+1$ and $b=p$ have none over the reals. A bare cycle is therefore not automatically a meaningful configuration recursion.

The source asks for a document-entry design, not a graph theorem. This investigation establishes a concrete bootstrap obstruction and a restricted loader rule; it does not choose or validate a new PreTeXt root element.
\section{Verification and adversarial audit}
The construction and counterexample are checked in the written finite model. No actual PreTeXt build, browser, EPUB reader, or publishing service was tested. All runtime and identity assumptions used by the proof are stated explicitly.
\paragraph{Audit corrections.} Cycle rejection is a declared one-time import policy. No assumption is made that actual PreTeXt loading behaves like the naive recursive example.
\paragraph{External computation.} None; symbolic arguments and boundary checks only.
\section{Scope and remaining gaps}
This is a restricted route investigation. The full feature requires implementation, supported-version decisions, conformance tests, and user-facing evaluation that were not performed here. The original agenda remains unresolved, and no novelty or completed software feature is claimed.
\paragraph{Integration review.} Codex (GPT-6; exact serving revision not exposed) reviewed the written arguments and expanded layout and source-reference wrappers. Repository metadata and the production web renderer were checked using Python 3.13 and JavaScript with MathJax 3.2.2. This model-assisted review is not independent human proof verification; no PDF compilation or formal proof checking is asserted.

\section{Completion Estimate}
\noindent\textbf{COMPLETION: 0\%}
The calculations below the preserved source statement investigate a concrete route,
but establish no substantiated new progress toward the entire original target.
Known special cases and failed sufficient conditions are not credited as a solution.
\section{References}
The original source references and catalogue attribution are preserved in the
Sources section above. The following URLs were actually consulted for this attempt;
retention of an additional original reference does not claim fresh consultation of it.
\begin{itemize}
\item \url{https://github.com/PreTeXtBook/community/wiki/Moderated-Problem-Session} (consulted 2026-10-11). 
\end{itemize}
\section{Model and contributor disclosure}
OpenAI gpt-6-astra, requested reasoning configuration ultra, generated this research
attempt on 2026-10-11. The exact serving revision is not exposed. Codex GPT-6
performed orchestration and repository-format review; its exact serving revision
is also not exposed. The mathematical route, boundary checks and final adversarial
self-audit followed contribution prompts A then B in
\texttt{docs/contribution\_prompts.md}, after reading \texttt{CONTRIBUTING.md}.
Python 3.14.2 assembled the standalone documents; any mathematical computation is
specified in Verification. No theorem prover was used. No human mathematical
authorship, substantial human proof correction, or independent human proof checking
is asserted. Model-assisted review, source retrieval, and structural checks do not
establish correctness or novelty. Full private reasoning transcripts are unavailable
as reproducibility artifacts; the complete arguments, inputs and outcomes actually
used are stated here or in the companion batch audit.
\section{Final status}
Unresolved as a new research contribution. This is a concrete route investigation,
not a claim of new complete resolution or a current-status certification. A known
answer reproduced for context does not become a new contribution. The catalogue's
source-status label is preserved; this attempt status does not assert that a
source-labelled solved question is open.
\end{document}
